\documentclass{article}
\usepackage{iclr2027_conference,times}
\input{math_commands.tex}
\usepackage{hyperref}
\usepackage{url}
\usepackage{booktabs}
\usepackage{multirow}
\usepackage{graphicx}
\usepackage{amsmath,amssymb}
\usepackage{xcolor}
\usepackage{tikz}
\usetikzlibrary{arrows.meta,positioning,shapes.geometric,fit}
\usepackage{enumitem}
\setlist{nosep,leftmargin=*}

% ---- Drafting helpers: set \showtodostrue -> \showtodosfalse for the camera-ready PDF ----
\newif\ifshowtodos \showtodostrue
% 2Wiki results: keep \twowikifalse unless the team completes 2Wiki for all four approaches
\newif\iftwowiki \twowikifalse
\newcommand{\todo}[1]{\ifshowtodos\textcolor{red}{\textsf{\small[TODO: #1]}}\fi}
\newcommand{\owner}[1]{\ifshowtodos\textcolor{blue}{\textsf{\scriptsize[Owner: #1]}}\fi}
% ---- Review highlights (Andy's purple notes): set \showupdatestrue -> \showupdatesfalse to hide them ----
\newif\ifshowupdates \showupdatestrue

\newcommand{\W}{\textsc{W}}
\newcommand{\Sg}{\textsc{S}}
\newcommand{\Hm}{\textsc{H}}
\IfFileExists{dsfont.sty}{\usepackage{dsfont}\newcommand{\ind}{\mathds{1}}}{\newcommand{\ind}{\mathbf{1}}} % indicator 1[.]

\title{RightRoute: Stage-Level Model Routing that\\Optimises Accuracy, Cost and Risk\\per Instance}

\author{Jonathan Dohyung Ahn, Andy Hock Lin Tai, Mohamed Dhameem \& Ziming Cheng\\
School of Computing and Information Systems, Singapore Management University
}

\iclrfinalcopy
\begin{document}
\maketitle

% =====================================================================
\begin{abstract}
Deployed AI is increasingly limited not by what models can do but by which model does each task. Most business workflows run in several stages (sub-tasks), and each stage needs a right-sized model: a small model has often already completed a stage correctly, while a costly upgrade may add nothing or even introduce an error. Yet existing routers choose one model per query, and compound-system optimisers fix one combination of models per dataset, so neither decides per instance and per stage. We study this decision under accuracy, cost and risk. Preliminary experiments gave three insights: (i) the strongest model is not always best; (ii) a router learns only from a carefully chosen view of the input, not the raw input; and (iii) an upstream choice changes downstream cost and accuracy. Building on them, we propose RightRoute: at each stage a router deliberately selects the model, keeping a cheaper model's result where it suffices and escalating to a stronger model where it does not or, for high-risk tasks, to a human, trained on a single reward that combines accuracy, cost and risk. Across two benchmarks and four families of routing and model-switching approaches, RightRoute (1) matches the accuracy of always using the strongest model at substantially lower cost, and exceeds it where the cheaper model succeeds on instances the strongest model fails; (2) escalates high-risk tasks to a human when a model error would cost more than the human; and (3) shows which approach suits which objective: fine-tuned executors for accuracy, selective routing for compute savings, and lightweight decision models for deployability.
\end{abstract}

% =====================================================================
\section{Introduction}

Organisations have invested heavily in generative AI, yet most see no measurable return, according to MIT's Project NANDA \citep{nanda2025genai}. The organisations that do capture value are distinguished less by the models they use than by how they redesign their workflows around them \citep{mckinsey2025state}, and analysts expect small, task-specific models to be used far more widely than general-purpose ones \citep{gartner2025small}. Redesigning a workflow breaks it into stages (sub-tasks), and a growing range of model sizes gives each stage several candidate models. The bottleneck in deployed AI is therefore no longer model capability but model allocation: choosing which right-sized model should execute each stage, for each request.

Not every stage needs the strongest model \citep{mckinsey2026agents}, but every stage needs the right one. For a routine request, a small model may read the request and take the action correctly, with a large model needed only to look up the right records; for an unusual, high-risk request, the large model may be needed up front and the final action may belong to a human. A costly upgrade can add nothing, or even turn a correct answer into a wrong one, and the cost of a mistake differs by stage: a wrong lookup is easy to detect and undo, a wrong medication order is not. Our preliminary results show that this matters in practice: choosing per instance and per stage could beat always using the strongest model at about a third of its cost (Figure~\ref{fig:headroom}, Appendix~\ref{app:prelim}). We therefore ask: \emph{can we choose the right-sized model for each stage of each request, to achieve similar or better accuracy at lower cost while accounting for risk in high-risk tasks?}

Existing work answers coarser questions. Query-level routers such as RouteLLM \citep{ong2025routellm} and Router-R1 \citep{zhang2025routerr1} choose one model for the whole request. Compound-system optimisers such as AgentTTS \citep{wang2025agenttts} and LLMSelector \citep{chen2025llmselector} choose a model per stage but fix that choice for every instance. Step-level routers such as TRIM \citep{trim2026} decide per instance, but over the repeated steps of a single reasoning chain, evaluated on math reasoning, rather than distinct stages. None places instance-specific risk in the objective. Making the decision per instance and per stage raises three challenges: (i) the router must decide before the Strong model runs, from limited information; (ii) stages are interdependent, so an upstream choice changes the cost and accuracy of the stages after it; and (iii) the cost of an error varies by instance, so the same decision can be acceptable for one request and unacceptable for another.

To address these challenges, we first ran preliminary experiments, which yielded three insights (\S\ref{sec:insights}): always using the strongest model is not the best policy; a router can learn only when its input is deliberately selected; and upstream choices affect downstream cost and accuracy. Guided by these insights, we propose RightRoute. It runs the Weak model first and, at each stage, learns from a deliberately selected view whether to keep the Weak result or escalate to a Strong model, or, for high-risk instances, to a human. It is trained on a single reward that combines accuracy, cost and a risk term that activates above a rubric-derived risk gate (\S\ref{sec:method}). Our primary method is an LLM router; we compare it with three alternative families---non-LLM routers, System One judges and fine-tuned executors with learned escalation---and evaluate them on HotpotQA and MedAgentBench against always using the strongest model (\S\ref{sec:experiments}--\ref{sec:results}).

Our experiments support three contributions:
\begin{enumerate}
\item \textbf{Right-sized models per stage reach or beat the always-Strong bar at substantially lower cost}, with a router trained from the outcome reward alone, without routing labels.
\item \textbf{Risk-aware routing with human escalation:} for high-risk tasks such as medication orders, a model error can cost more than a human, so the reward includes a risk term and RightRoute escalates to a human when even the strongest model is likely to fail and the expected cost of that error exceeds the cost of a human.
\item \textbf{The four approaches reduce to three deployment choices, one per objective:} fine-tuned executors for accuracy, selective routing for compute savings, and lightweight decision models for deployability.
\end{enumerate}

% =====================================================================
\section{Related Work}
\label{sec:related}

We organise prior work by \emph{routing unit} (whole query, repeated steps, or distinct stages) and by whether the decision is \emph{learned per instance} or \emph{fixed per dataset}; Table~\ref{tab:gap} (Appendix~\ref{app:related}) gives the full map.

\textbf{Whole query.} RouteLLM \citep{ong2025routellm} learns a router from preference data, AutoMix \citep{aggarwal2024automix} cascades using self-verification, and Router-R1 \citep{zhang2025routerr1} trains an LLM router with reinforcement learning, including on HotpotQA \citep[see also][]{chen2023frugalgpt,ding2024hybrid,dekoninck2025unified}; conformal LLM routing \citep{uddin2026conformal} adds a dataset-level guarantee from the question alone. All choose one model for the whole request.

\textbf{Repeated steps.} TRIM \citep{trim2026} escalates single steps of a math reasoning chain after the cheap model writes them, scored by a step grader, with no risk term; R2-Reasoner \citep{r2reasoner2025} and Division-of-Thoughts \citep{shao2025dot} route sub-tasks that a decomposer invents per query; RSD \citep{liao2025rsd} and SpecReason \citep{pan2025specreason} accept cheap steps by a fixed threshold rule.

\textbf{Distinct stages.} AgentTTS \citep{wang2025agenttts}, LLMSelector \citep{chen2025llmselector} and Cognify \citep{he2025cognify} allocate models per stage or module but return one setting per dataset. MasRouter \citep{yue2025masrouter} routes inside multi-agent systems without reading stage outputs, and TRAQ \citep{li2024traq} builds per-stage conformal sets without choosing models. Learned escalation in multi-stage architectures is described as largely unexplored \citep{routingsurvey2026}: no prior system combines fixed heterogeneous stages, a learned per-instance decision, reading the Weak stage's output before the Strong model is spent, and instance-specific risk in the decision.

\textbf{Escalation signals and decision models.} Cascades usually escalate on the Weak model's own signals, such as token probabilities, self-verification \citep{chen2023frugalgpt,aggarwal2024automix}, P(True) \citep{kadavath2022know} or semantic entropy \citep{farquhar2024semantic}, which need access to the Weak model or extra calls; Self-RAG and Adaptive-RAG decide when to retrieve \citep{asai2024selfrag,jeong2024adaptiverag}. Small decision models that answer typed questions over text, such as nano-jev \citep{nanojev2026} and the Strands Decider \citep{strandsdecider2026}, are offered as agent gates; we found no evaluation of their pretraining as a router's starting point, or of routing from an untrained typed question.

% =====================================================================
\section{Problem Definition and Preliminary Insights}
\label{sec:insights}

\subsection{Problem definition}
\label{sec:problem}

Let $i$ index \emph{instances}, where an instance $x_i$ is one task (e.g., one HotpotQA question or one MedAgentBench EHR request) decomposed into $K$ \emph{stages} (sub-tasks), $k=1,\dots,K$ ($K$ may vary by instance). At each stage a policy assigns an executor $a_{i,k}$: the Weak model \W{}, small and cheap; the Strong model \Sg{}, large and costly; or, for high-risk instances, a human \Hm{} (\emph{human escalation}). The action set $\mathcal{A}_i$ is $\{\W,\Sg\}$ when the instance's risk score $\tau_i\in[0,1]$ is below the risk gate $\tau_{\text{gate}}=0.8$, and $\{\W,\Sg,\Hm\}$ otherwise. Each stage's output $o_{i,k}$ is passed to the next, so the decision at stage $k$ can use only $s_{i,k}=(x_i;\, o_i^{(<k)};\, a_i^{(<k)})$: the task and, for every earlier stage, its output and executor (notation: Table~\ref{tab:notation}, Appendix~\ref{app:notation}). Decisions are made per stage, but accuracy, cost and risk are optimised per instance: the goal is a policy $\pi(a_{i,k}\mid s_{i,k})$ that maximises the expected instance reward $R_i$ (\S\ref{sec:reward}). The \emph{oracle} is the best choice per instance in hindsight (a correct action sequence if one exists, the cheapest when several are correct); it upper-bounds any policy over the given models. In HotpotQA, $K{=}2$: retrieval selects the supporting passages from 100 candidates, and answering reads them and answers, giving four fixed policies \W\W, \W\Sg, \Sg\W{} and \Sg\Sg{} (retrieval choice first). In MedAgentBench, each EHR task is a multi-turn, tool-using agent episode whose rounds are the stages ($K\le8$): at each round the executor reads the patient record, calls a tool or returns the result, and a decision may be taken at any round to keep the current executor or to re-run the task with another (\Sg, or \Hm{} when $\tau_i\ge\tau_{\text{gate}}$).

\textbf{Risk score.} $\tau_i$ comes from a three-item rubric scored 1--5 (Appendix~\ref{app:rubric}): impact of error ($Q_1$), detectability ($Q_2$) and reversibility ($Q_3$). Then $\tau_i = (15 - (Q_1+Q_2+Q_3))/12$, with a floor of 0.8 when $Q_1{=}1$ and clipping to $[0.10, 0.95]$, so a higher $\tau_i$ means a riskier instance. $Q_1$ and $Q_2$ build on the severity and detection concepts of FMEA \citep{iec60812}, rescaled to 1--5 with 5 as lowest risk; $Q_3$ is our addition for routing, because two stages with equal impact and detectability can need different models depending on whether an error can be undone. The rubric is frozen before training.

\textbf{Cost.} Cost is task-model compute as a share of always-Strong on the same instances (\S\ref{sec:experiments}); the decision's own compute is reported separately. The preliminary analysis also uses token-aware prices in which answering cost depends on the retrieval choice, because the answering prompt carries the selected evidence (Table~\ref{tab:passcost}).

\subsection{Preliminary experimental findings}
Before the main experiments, we ran a preliminary study on HotpotQA with two stages: retrieval from 100 candidate passages, then answering. Qwen2.5-1.5B-Instruct is the Weak model and Qwen2.5-7B-Instruct the Strong model; all four Weak/Strong combinations were run on the 7,405 questions, and results are reported on the 1,481 held-out ones. Our first attempt failed: neither a trained router nor six feature families could predict which model wins where the two disagree (AUC 0.51--0.57, near chance), although question difficulty was predictable (AUC 0.61--0.69). Comparing the four outcomes question by question instead---where they agree, what any of them gets right, and where Weak succeeds but Strong fails---yielded the three insights below (Appendix~\ref{app:prelim}).

\subsection{Three insights on per-instance, per-stage routing}
\textbf{Insight 1: Strong is not always best.} On one question in eleven, a cheaper configuration succeeds where always-Strong fails, and the oracle reaches 68.6\% against 59.6\% at 32\% of the cost (Figure~\ref{fig:headroom}, Appendix~\ref{app:prelim}). AgentTTS finds a preferred model per sub-task and dataset; we find the preference also changes from question to question.

\textbf{Insight 2: input decides learnability.} Shown the raw input (the question and all 100 passages), our first routers collapsed to a constant policy. Two changes made the decision learnable: show the router the passages the Weak retriever selected, and ask whether they are enough---a question that always has an answer, unlike ``which model wins'', which is undefined on about four questions in ten. The same architecture then ranks well (AUC ${\approx}0.77$ on validation); on MedAgentBench the view is the Weak agent's attempt. AgentTTS has no counterpart: it searches configurations rather than learning a decision.

\textbf{Insight 3: upstream choice sets downstream accuracy and cost.} Weak retrieval hands the answerer about 850 tokens and Strong retrieval 330, so the same answering call costs $2.6\times$ more on one branch; and the Weak answerer is right 48\% of the time on Strong-retrieved evidence but 38\% on Weak-retrieved. This is the per-instance form of AgentTTS's third insight, and it is why RightRoute prices each stage by its branch and decides the stages in sequence.

% =====================================================================
\section{RightRoute: Method}
\label{sec:method}

RightRoute turns the three insights into three design choices. \textbf{Strong is not always best} (Insight~1): since the right executor changes from instance to instance and stage to stage, the router is trained on an outcome reward that prices accuracy, cost and risk, rather than to imitate the strongest model. \textbf{Input decides learnability} (Insight~2): the router never reads the raw input, only the Weak model's own attempt at that stage, so the Weak model always goes first, and each stage has its own router asking one question: is this attempt enough? \textbf{Upstream choice sets downstream accuracy and cost} (Insight~3): each stage's output is the next stage's input, so the stages are decided in sequence and every cost in the reward is conditional on the branch taken.

\subsection{Weak-first escalation}
\label{sec:weakfirst}
At each stage the Weak model goes first and produces an attempt $w_{i,k}$. The router then decides from everything known so far, $s_{i,k}=(x_i;\, o_i^{(<k)};\, a_i^{(<k)};\, w_{i,k})$---the instance, the earlier stages' outputs and executors, and the attempt itself---whether to keep the Weak result (\W), escalate the stage to the Strong model (\Sg) or, when the instance is above the risk gate, hand the task to a human (\Hm). The router never writes the stage output; whichever result is kept or produced by the escalation passes to the next stage, and a human completes the task. Strong compute is therefore spent only where the router asks for it, but the attempt it replaces is not free: an escalation costs the Weak attempt plus the Strong call, on both benchmarks. On HotpotQA the Strong-everywhere path costs 1,009 units against 929 for the Strong model alone; on MedAgentBench an escalation re-runs the whole task after the Weak attempt. An escalated stage thus always costs slightly more than always-Strong at that stage, and the budget term of Eq.~\ref{eq:reward} acts as a small, constant escalation penalty: the price of looking before escalating.

\subsection{Reward}
\label{sec:reward}
The router is trained on the outcome, not by imitating the strongest model:
\begin{equation}
R_i = Q_i - \lambda\,\frac{C_i}{C_{SS,i}} - \max\!\Big(0, \frac{C_i}{C_{SS,i}} - 1\Big) - \lambda_t\Big[\frac{C_{F,i}}{C_{H,i}}\,\ind[\text{fail}_i]\,\ind[a_i\neq \Hm] + \ind[a_i = \Hm]\Big]\,\ind[\tau_i \ge \tau_{\text{gate}}].
\label{eq:reward}
\end{equation}
Here $Q_i$ is answer quality (token-F1 on HotpotQA; the official pass/fail grader on MedAgentBench; 1 when a human completes the task); $C_i$, $C_{SS,i}$, $C_{H,i}$ and $C_{F,i}$ are the compute of the chosen path, the always-Strong compute of the same instance, the human cost and the failure cost (rework plus slip-through $\times$ harm). The second term applies steady cost pressure ($\lambda{=}0.10$); the third penalises spending beyond the instance's always-Strong budget, which avoids false penalties on long agentic tasks. The risk term applies only at or above the risk gate (MedAgentBench's medication orders; no HotpotQA instance). There, a model failure costs $C_F/C_H$ and a human costs one ($\lambda_t{=}1$), so the reward escalates to a human when the expected failure cost exceeds the human cost. The reward is used only in training; at inference the router sees neither the reward nor the ground truth.

\subsection{Primary method and three alternatives}
\label{sec:families}

\textbf{Primary method: LLM router.} One router per stage, a Qwen2.5-1.5B-Instruct model with a LoRA adapter \citep{hu2022lora} (rank 16; 4.4M trained parameters), reads the Weak model's attempt (Insight~2) and gives $P(\W)$, the probability that the Weak result is enough. The Weak result is kept when $P(\W)$ clears a threshold $\theta$ set on validation and is escalated otherwise (Figure~\ref{fig:framework}). The answering router also sees which retriever ran and what it selected, so the second decision depends on the first (Insight~3). On MedAgentBench one decision is taken per task, after the Weak agent's full attempt or at its second round. We train the router three ways, changing one thing at a time. \emph{SFT (label)} learns from per-stage labels: for retrieval, \W{} if the Weak retriever's passages contain the gold answer, \Sg{} if only the union of both retrievers' passages does, and X if neither (trained as \W, because escalation cannot fix it); for answering, whether a checker finds the Weak answer supported by its evidence. \emph{SFT (label) + Reward} continues from those weights on the expected reward of Eq.~\ref{eq:reward}. \emph{GRPO(Reward)} starts from a fresh adapter and learns from Eq.~\ref{eq:reward} alone, with no labels, by group-relative policy optimisation \citep{shao2024deepseekmath} ($\beta{=}0.03$), so it carries over to settings where no labels exist; the both-fail case needs no rule, because escalating adds cost and leaves $Q$ unchanged.

\textbf{Three alternatives.} We compare three alternatives under the same splits, reward and read-once protocol (full descriptions in Appendix~\ref{app:families}). The \emph{non-LLM router} makes the same per-stage choice with a lightweight classifier (linear or MLP; on MedAgentBench also random forest, SVM and logistic regression) over lexical, structural and stage-aware features, testing whether an LLM is needed at all. The \emph{System One judge} asks \emph{will the Weak result do?} with a small decision model that scores options without generating: nano-jev \citep{nanojev2026}, a 33.4M-parameter MiniLM cross-encoder \citep{wang2020minilm}, and the 2B and E4B Strands Deciders \citep{strandsdecider2026}, trained on the same per-stage labels as SFT (label) on HotpotQA. The \emph{fine-tuned executors with learned escalation} replace routing between fixed models: we fine-tune the retrievers (Qwen2.5 7B--72B) and answerers (Llama 3B--70B) on HotpotQA's training set, run the smallest models first and let fixed or LLM-proposed rules decide when to add a larger one, with a human as the last rung on high-risk MedAgentBench tasks.

% Andy's 22 Nov update, in purple; hidden by the same \showupdatesfalse
\definecolor{andycolor}{RGB}{128,0,128}\DeclareRobustCommand{\andy}[1]{\ifshowupdates\textcolor{andycolor}{#1}\else#1\fi}\DeclareRobustCommand{\andynote}[1]{\ifshowupdates\textcolor{andycolor}{\textsf{\scriptsize[Update (Andy): #1]}}\fi}%

\begin{figure}[t]
\centering
\includegraphics[width=\linewidth]{figs/routing.png}
\caption{Weak-first stage-level routing. At each stage the Weak model runs first; the router reads the instance, the upstream output and the Weak attempt and decides, before any Strong cost is spent, whether to keep the Weak result or escalate to the Strong model or, for high-risk instances ($\tau_i \ge 0.8$), to a human, who completes the stage (in our MedAgentBench runs, the task). Each later stage repeats the pattern on the previous stage's output. HotpotQA has two stages; on MedAgentBench one decision is taken per task (\S\ref{sec:families}). The reward is computed once per instance (Eq.~\ref{eq:reward}).}
\label{fig:framework}
\end{figure}

% =====================================================================
\section{Experimental Setup}
\label{sec:experiments}

\textbf{Datasets.} \emph{HotpotQA} \citep{yang2018hotpotqa} supplies 7,405 multi-hop questions, each with the 100 candidate passages distributed by AgentTTS. We split them 4,814 / 1,110 / 1,481 (train / validation / held-out test), frozen before any training. Our fine-tuned models learn from HotpotQA's own training set rather than our 4,814 questions, but are tuned on the same 1,110 validation questions and tested once on the same 1,481 held-out test questions as every other method. \emph{MedAgentBench} \citep{jiang2025medagentbench} supplies 300 EHR agent tasks run against a live FHIR server, split 90 / 60 / 150 with no patient shared across splits. The server runs in the benchmark's own container on UTC, so tasks that ask about a recent time window are scored as the benchmark intends. The two medication-order task types score $\tau\ge0.8$ under the rubric (30 of the 150 held-out test tasks) and are the only instances where human escalation is an option. Every HotpotQA instance falls below the risk gate, so the risk term is inactive there.

\textbf{Executors.} On HotpotQA the Weak/Strong pairs are AgentTTS's own \citep{qwen2025qwen25,grattafiori2024llama3}: Qwen2.5-7B / Qwen2.5-72B for retrieval and Llama-3.2-3B / Llama-3.1-70B for answering, full precision, greedy decoding. On MedAgentBench they are Qwen2.5-7B and Qwen2.5-72B-Instruct-AWQ, served locally, with human escalation on high-risk tasks; the published MedAgentBench figure was obtained on a since-retired hosted FP8 serving (Together), and our local AWQ serving of the same model scores 56.7\% on all 300 tasks and 54.0\% on the 150 held-out test tasks. The preliminary study of \S\ref{sec:insights} used Qwen2.5-1.5B / 7B at both stages.

\textbf{Metrics.} Accuracy on HotpotQA is lenient containment of the reference answer, the rule the AgentTTS authors confirmed; under it, our reproduction of AgentTTS scores 74.6\% on a 500-question slice of the same dev set (theirs is unreleased), against their published 74.0\%. On MedAgentBench it is the official grader's pass rate. Cost is task-model compute as a share of always-Strong on the same instances, in AgentTTS's normalised units (Table~5 of \citealp{wang2025agenttts}) on HotpotQA and normalised FLOPs on MedAgentBench. The router's own compute, about 3.6 points of always-Strong on HotpotQA and 19\,ms per decision on MedAgentBench, is reported separately in Appendix~\ref{app:metrics} with the supporting metrics.

\textbf{Baselines.} Always-Strong, always-Weak and every fixed per-stage policy; AgentTTS, published and reproduced on our pipeline; the published MedAgentBench figure; random escalation at matched volume; for the fine-tuned executors, the best single fine-tuned retriever--answerer pair and the fine-tuned 7B alone on MedAgentBench; and entropy and logistic-regression cascades for the judges. \andynote{the best pair was picked by looking at test scores; it should be picked on the validation questions first and then tested once (to do before submission)}

\textbf{Protocol and statistics.} Every experiment was registered with its predictions before it ran; thresholds and checkpoints are chosen on validation, and each trained configuration is tested exactly once on the held-out test. All methods are scored on the same test instances; single-policy rows carry 95\% Wilson intervals (about $\pm2.2$ points on HotpotQA, $\pm7.7$ on MedAgentBench). Routers are single-seed unless stated; this is a stated limitation.

% =====================================================================
\section{Results and Analysis}
\label{sec:results}



\begin{table}[t]
\caption{\textbf{Best result per method on the held-out tests} (highest-accuracy configuration per method; Figure~\ref{fig:results} plots accuracy against cost). Cost = task-model compute as \% of the Strong model alone on the same instances. $^\ast$With human escalation on high-risk tasks (a human counted correct). Notes $\ddagger$, $\S$ and $\P$, the cost conversion, the fine-tuned configurations and the routers' and judges' own compute are in Appendix~\ref{app:metrics}. Single runs unless noted.}
\label{tab:main}
\centering\small
\resizebox{\linewidth}{!}{%
\begin{tabular}{@{}lllcccc@{}}
\toprule
& & & \multicolumn{2}{c}{HotpotQA ($n{=}1{,}481$)} & \multicolumn{2}{c}{MedAgentBench ($n{=}150$)}\\
\cmidrule(lr){4-5}\cmidrule(lr){6-7}
Method & Training & Trained params & Acc.\ (\%) & Cost (\%) & Acc.\ (\%) & Cost (\%)\\
\midrule
Always-Strong & -- & -- & 76.8 & 100 & 54.0 & 100\\
Always-Weak & -- & -- & 50.3 & 8.6$^\ddagger$ & 12.0 & ${\approx}$23\\
Oracle (hindsight upper bound) & -- & -- & 81.5 & 23.3 & 63.3 & 54.0\\
Published (AgentTTS / MedAgentBench) & -- & -- & 74.0 & ${\approx}$101 & 51.3$^\S$ & --\\
AgentTTS (reproduced, 500-question slice) & search & -- & 74.6 & ${\approx}$101 & -- & --\\
\midrule
LLM router & SFT (label) & 4.4M & 76.7 & 89.9 & -- & --\\
LLM router & GRPO(Reward), label-free & 4.4M & 75.9 & \textbf{69.7} & 61.3 & 66.6\\
LLM router + human escalation & GRPO(Reward) + priced risk & 4.4M & -- & -- & 77.3$^\ast$ & 76.0 + SGD 184$^\P$\\
Non-LLM router & SFT (label) & 3.1M$\mid$6K & 76.8 & 98.2 & 63.3 & 95.6\\
System One judge (E4B$\mid$nano-jev) & SFT (label) & 109M$\mid$33M & 76.7 & 79.4 & 63.3 & 72.6\\
Fine-tuned executors & SFT & 7B--70B & \textbf{81.3} & 21 & 76.7 & 9.3\\
Fine-tuned executors + escalation (larger models, then human) & SFT + learnt rules & 7B--72B & 79.6 & -- & \textbf{97.3}$^\ast$ & 89.8\\
\bottomrule
\end{tabular}}
\end{table}

\begin{figure}[t]
\centering
\includegraphics[width=0.95\linewidth]{figs/results.pdf}
\caption{Accuracy against cost for the best configuration of each method (Table~\ref{tab:main}). Upper-left is better. The oracle is a hindsight upper bound for routing between the two untrained models; fine-tuned executors are not bound by it. $^\ast$Includes human escalation on high-risk tasks (counted as correct); the LLM router + human point shows model cost only, plus SGD 184 of human time. Costs are \% of the Strong model alone.}
\label{fig:results}
\end{figure}



\textbf{Headroom exists, and part of it is invisible to query-level routing.} Insight~1 shows that per-stage routing is worth attempting, not that a learned router captures the gain; the remaining results test that. On HotpotQA with the main executors, restricting the oracle to one model for the whole query (the better of \W\W{} and \Sg\Sg{} per question) gives 80.1\% at 38.4\% of always-Strong; allowing the choice per stage gives 81.5\% at 23.3\%. The difference, 1.4 points of accuracy and 15 points of cost, exists only at stage granularity: mixed paths answer some questions that neither single-model path does, and answer others more cheaply. This is the part a query-level router such as RouteLLM cannot reach. (The oracle selects nothing, so computing it on the held-out test does not break the read-once protocol.)

\textbf{Headroom on MedAgentBench.} With Qwen2.5-7B / 72B, the per-instance oracle reaches 63.3\% at 54.0\% of the Strong agent's cost, against 54.0\% for always-Strong ($+9.3$ points) and 12.0\% for always-Weak (at about 23\%); 36.7\% of tasks fail under both models, so no Weak/Strong policy can exceed it.

\textbf{Routers reach always-Strong accuracy at lower cost.} On HotpotQA every router family lands within one point of always-Strong (76.8\%; 95\% interval about $\pm2.2$ points), so the differences are in cost: 69.7\% for the label-free GRPO router, 79.4\% for the System One E4B judge, 89.9\% for SFT (label) and 98.2\% for the non-LLM classifier, which almost always chooses Strong. On MedAgentBench the GRPO router reaches 61.3\% at 66.6\% of the cost, capturing 79\% of the gap to the oracle, and the System One judge 63.3\% at 72.6\%, 7--9 points above always-Strong (54.0\%); the non-LLM router and the second-round-view router also reach 63.3\%, at nearly always-Strong cost; with 150 tasks (interval about $\pm7.7$ points) we read this as a consistent gain, not a proven one. A cheaper path can be \emph{more} accurate because the Strong model does not dominate the Weak one on every instance.

\textbf{Ablations tied to the insights.} Two design choices are tested directly: replacing the reward by labels alone (Insight~1) raises cost from 69.7\% to 89.9\% at similar accuracy (Table~\ref{tab:ablation}), and a router reading the raw input instead of the Weak model's attempt (Insight~2) collapses to a constant policy (Appendix~\ref{app:negative}).

\textbf{Training recipe ablation (LLM router).} Changing one element at a time (Table~\ref{tab:ablation}), the reward pass cuts cost by 18 points for $-0.7$ points of accuracy, and removing the labels keeps the same operating point: the outcome reward alone suffices, and this is the only router of the four approaches trained without routing labels. On MedAgentBench, which has no per-stage labels, GRPO is the only option; the second-round view reaches the oracle's accuracy at always-Strong cost.

\begin{table}[!htb]
\caption{\textbf{LLM router: training-recipe ablation} (Qwen2.5-1.5B base; held-out test; $\theta$ chosen on validation). Labels exist only on HotpotQA; the risk term is inactive there (uniform low $\tau$). Cost and baselines as in Table~\ref{tab:main}; risk-term variants: Table~\ref{tab:risk}.}
\label{tab:ablation}
\centering\small
\begin{tabular}{@{}lcccc@{}}
\toprule
& \multicolumn{2}{c}{HotpotQA} & \multicolumn{2}{c}{MedAgentBench}\\
\cmidrule(lr){2-3}\cmidrule(lr){4-5}
Recipe (one change per row) & Acc.\ (\%) & Cost (\%) & Acc.\ (\%) & Cost (\%)\\
\midrule
Always-Strong & 76.8 & 100 & 54.0 & 100\\
\midrule
SFT (label) & 76.7 & 89.9 & -- & --\\
SFT (label) + Reward & 75.96 & 71.6 & -- & --\\
GRPO(Reward), label-free (full-attempt view) & 75.89 & 69.7 & 61.3 & 66.6\\
GRPO(Reward), label-free (second-round view) & -- & -- & 63.3 & 100.2\\
\bottomrule
\end{tabular}
\end{table}

\textbf{Router size ablation (LLM router).} Varying one router's size from 0.5B to 7B (Table~\ref{tab:routersize}; same SFT (label) recipe, other router fixed at 1.5B, validation) leaves accuracy flat, within 3 questions, so cost decides. Larger routers rank better but cost more: a 7B retrieval router costs more than always-Strong once both router calls are charged, while the 1.5B pair is cheapest at both stages, which the full 16-pairing grid confirms (Table~\ref{tab:size}, Appendix~\ref{app:extra}). An untrained router ranks at chance (AUC 0.466) and an 8-feature classifier reaches 0.728, below every trained router. Routers should be selected on composed cost at equal accuracy, not on AUC.

\begin{table}[!htb]
\caption{\textbf{LLM router size ablation} (validation, $n{=}1{,}110$; one stage's router varied, the other fixed at 1.5B). Accuracy is flat across sizes (76.76--77.03\%, always-Strong 76.76\%); cost (\% of always-Strong, 929 units) charges both router calls.}
\label{tab:routersize}
\centering\small
\begin{tabular}{@{}lccccc@{}}
\toprule
& \multicolumn{2}{c}{Retrieval router} & \multicolumn{2}{c}{Answering router} & \\
\cmidrule(lr){2-3}\cmidrule(lr){4-5}
Size & AUC & Cost (\%) & AUC & Cost (\%) & Adapter\\
\midrule
0.5B & 0.851 & 84.6 & 0.670 & 92.5 & 2.2M\\
1.5B & 0.870 & \textbf{74.3} & 0.798 & \textbf{74.3} & 4.4M\\
3B & 0.886 & 94.2 & 0.853 & 95.4 & 7.4M\\
7B & 0.897 & 116.2 & 0.876 & 100.2 & 10.1M\\
\bottomrule
\end{tabular}
\end{table}

\textbf{Risk tier: the priced reward escalates every high-risk order to a human.} On the 30 high-risk MedAgentBench test orders (Table~\ref{tab:risk}), the router without the risk term escalates none and the model fails on 21; adding the risk term without prices escalates 27 and leaves 2 failures; pricing human time and failure cost (Eq.~\ref{eq:reward}) escalates all 30 and leaves none, at a model cost of 76.0\% of always-Strong plus SGD 183.75 of human time, whereas always-Strong's 21 unseen failures would cost SGD 13,863 at the same rates. The priced router is conservative (9 of its 30 escalations were not strictly needed) and identical across three seeds.

\textbf{Fine-tuning beats scale; learned switching does not beat the best fixed pair.} Fine-tuned executors give the highest accuracy: on HotpotQA a fine-tuned 7B retriever with a 70B answerer matches the two-model oracle's accuracy (81.3\% vs.\ 81.5\%) at 21\% of always-Strong's cost, and a 0.8B$\to$8B pair reaches 80.8\% at 1.4\% (Figure~\ref{fig:ftexec}); on MedAgentBench a fine-tuned 7B alone (76.7\%) exceeds the oracle, and escalating up to a human reaches 97.3\%. Learned escalation rules never beat the best fixed fine-tuned pair (Appendix~\ref{app:extra}).

\begin{table}[!htb]
\caption{\textbf{High-risk tier on MedAgentBench} (30 medication-order tasks in the 150-task held-out test). Accuracy counts a human escalation as correct. Cost is \% of the Strong model alone, as in Table~\ref{tab:main}. Human cost = human escalations $\times$ time $\times$ rate; failure cost = $C_F$ per high-risk task completed wrongly by a model.}
\label{tab:risk}
\centering\small
\resizebox{\linewidth}{!}{%
\begin{tabular}{@{}lccccc@{}}
\toprule
& Acc.\ (\%) & Cost (\%, model + human) & Escalated to a human & Model failed & Mean $R$\\
\midrule
Always-Weak & 12.0 & 23.0 & 0 & 30 & 0.095\\
Always-Strong & 54.0 & 100 & 0 & 21 & 0.164\\
GRPO(Reward), no risk term (\W/\Sg) & 61.3 & 66.6 & 0 & 21 & 0.377\\
GRPO(Reward) + risk, unpriced (\W/\Sg/\Hm) & 76.0 & 116.2 (77.4 + 38.9) & 27 & 2 & 0.358\\
GRPO(Reward) + risk, priced (\W/\Sg/\Hm) & 77.3 & 76.0 + SGD 183.75 & 30 (21 needed) & 0 & 0.299\\
\bottomrule
\end{tabular}}
\end{table}

\textbf{Summary of the three claims.} Table~\ref{tab:claims} (Appendix~\ref{app:extra}) links each claim to its evidence and to what remains uncertain.



% =====================================================================
\section{Conclusion and Future Work}
We asked whether a workflow can choose the right-sized model for each stage of each request, achieving similar or better accuracy at lower cost while accounting for risk. Preliminary experiments showed that always using the strongest model leaves large headroom and yielded three insights that shaped RightRoute: a Weak-first router per stage that reads deliberately selected input and is trained on a single reward of accuracy, cost and risk.

Our held-out results support three conclusions. First, our primary method, the LLM router trained from the outcome reward alone without routing labels, reaches the always-Strong bar at lower cost on HotpotQA and exceeds it on MedAgentBench. Second, separate configurations for low- and high-risk tasks work: below the risk gate the router chooses between models, and above it, with failure and human costs in the reward, every high-risk MedAgentBench order was escalated to a human, with no model failures on that tier. Third, the four approaches reduce to three practical choices, each measured against always-Strong on the same instances: fine-tuned executors for accuracy, selective routing for compute savings, and lightweight decision models for deployability.

Three directions follow, each addressing a limitation in the Ethics Statement and Limitations. (i) \emph{More datasets and domains:} \iftwowiki TaskBench\else 2WikiMultiHopQA and TaskBench\fi{} for the routing methods, and GPAgentBench-2K \citep{chen2026gpagentbench} for the risk function, where escalation to a human specialist (abstain-and-refer) is a native action. (ii) \emph{More executor families and model sizes:} models beyond Qwen and Llama, with an ablation of which sizes are right for each dataset. (iii) \emph{A pilot in a real workflow:} measured human and failure prices would replace our assumed ones.

% =====================================================================
\section*{Ethics Statement and Limitations}
\textbf{Ethics.} RightRoute decides when a task is handed to a human, so its errors fall on people: under-escalation lets model mistakes through on high-stakes steps, and over-escalation adds reviewer workload and invites automation bias. The risk gate is meant to keep cheaper routing from lowering the safety bar, but its prices and rubric are our own and need clinical review before any clinical use, all the more because MedAgentBench (MIT licence, research use) is built from de-identified real patient cases from Stanford's STARR data warehouse, with patient-level timestamp jittering; we use it as distributed, add no patient data and run it only against its own FHIR container. HotpotQA is public (CC BY-SA 4.0), and no human subjects were involved. \textbf{Integrity and disclosures.} We report all results as obtained, including negative findings (Appendix~\ref{app:negative}), and read each held-out test once. The authors declare no competing interests; J.A.'s EngD is industry-embedded with an agentic AI start-up, which provided no funding, data or review of this work. Compute was provided by the SMU cluster and our lab's GPUs. Claude (Anthropic) assisted with drafting and editing text, writing analysis and training scripts, and preparing slides; the first author reviewed every prompt and line of code before any run, every experiment was pre-registered with predictions, and all reported numbers come from the stored outputs of those runs. The models studied are those named in \S\ref{sec:experiments}; no AI tool generated or edited experimental results.

\textbf{Limitations.} (1)~\emph{Assumed prices and rubric.} Clinician time, slip-through rate and incident cost come from the literature and our judgement, and the authors scored the risk rubric without an inter-rater check. The prices change the outcome: if a human's time costs the same as running the strongest model, the fine-tuned system sends all 30 high-risk test tasks to a human; at twice that cost, it sends 18 and two model mistakes go unnoticed. Direction~(iii) addresses this. (2)~\emph{Small evidence.} The risk results rest on 30 instances of two task types in one benchmark; most rows are single runs (the GRPO routers were identical across three seeds); the fine-tuned system's MedAgentBench escalation rules were learnt from 60 practice tasks; and MedAgentBench differences of 7--9 points are within the uncertainty of $n{=}150$. Direction~(i) addresses this. (3)~\emph{Scope.} Two benchmarks, two model families and greedy open-weight executors; wall-clock cost can differ from normalised units, and counting human escalation as correct measures the routing policy, not human accuracy. Directions~(i) and~(ii) address this. (4)~\emph{Evaluation details.} Three details could move our margins by up to about 3 points: overlap between the fine-tuned model's worked examples and its rule-learning tasks, a repetition setting on the MedAgentBench always-Strong baseline, and the order of passages in AgentTTS's HotpotQA files (Appendix~\ref{app:evaldetails}).

\section*{Reproducibility Statement}

Splits were frozen and committed before training, and each held-out test was read once per configuration; System One operating points and split hashes were committed before each test read, and the cross-dataset study was pre-registered. Code, configurations, run manifests and per-instance predictions are in the team repository (\url{https://github.com/mohameddhameem/STAIR}, commit \texttt{fc25f38}); the published judges, hyperparameters, prompts, rule learning and compute are in Appendices~\ref{app:impl}--\ref{app:compute}.

\bibliography{final_refs}
\bibliographystyle{iclr2027_conference}

% =====================================================================
\newpage
\appendix

\section{Notation}
\label{app:notation}
\begin{table}[h]
\caption{Notation used in the paper.}
\label{tab:notation}
\centering\small
\begin{tabular}{@{}lp{10.6cm}@{}}
\toprule
Symbol & Meaning\\
\midrule
$i$, $x_i$ & Instance index; instance $i$, the task itself (a HotpotQA question or a MedAgentBench EHR request)\\
$K$, $k$ & Number of stages (sub-tasks), which may vary by instance ($K{=}2$ on HotpotQA; $K \le 8$ rounds on MedAgentBench); current stage, $k=1,\dots,K$\\
\W, \Sg, \Hm & Weak model (small, cheap), Strong model (large, costly), human\\
$a_{i,k}$ & Executor chosen at stage $k$ (\W, \Sg{} or \Hm): the policy's decision\\
$o_{i,k}$ & Output of stage $k$, passed to the next stage\\
$o_i^{(<k)}$, $a_i^{(<k)}$ & Outputs of, and executors chosen at, all earlier stages $1,\dots,k-1$\\
$s_{i,k}$ & State before the decision at stage $k$: $(x_i;\, o_i^{(<k)};\, a_i^{(<k)})$\\
$w_{i,k}$ & Weak model's attempt at stage $k$ (RightRoute's weak-first design, \S\ref{sec:method})\\
$\mathcal{A}_i$ & Action set of instance $i$: $\{\W,\Sg\}$ below the risk gate, $\{\W,\Sg,\Hm\}$ at or above it\\
$\pi(a_{i,k}\mid s_{i,k})$ & Routing policy\\
$\tau_i$, $\tau_{\text{gate}}$ & Risk score of instance $i$ in $[0,1]$ (higher is riskier); risk gate, $0.8$\\
$Q_1, Q_2, Q_3$ & Rubric items, each 1--5 (5 = lowest risk): impact of error, detectability, reversibility\\
$R_i$ & Instance reward (Eq.~\ref{eq:reward})\\
$Q_i$ & Answer quality of instance $i$ (token-F1 on HotpotQA; official pass/fail on MedAgentBench)\\
$C_i$, $C_{SS,i}$ & Cost of the chosen action path; always-Strong cost of the same instance\\
$C_{H,i}$, $C_{F,i}$ & Human cost (time $\times$ rate); failure cost (rework plus slip-through $\times$ harm)\\
$\lambda$, $\lambda_t$ & Weight of the cost term; weight of the risk term\\
$\theta$ & Decision threshold on the router's score, chosen on validation\\
\bottomrule
\end{tabular}
\end{table}

\section{Related-work mapping}
\label{app:related}
\begin{table}[h]
\caption{Prior work by routing unit and decision type. The still-empty cell is learned, per-instance routing across distinct stages, where each stage's decision is made from the Weak attempt's output before the Strong model is spent, with instance-specific risk in the decision.}
\label{tab:gap}
\centering\small
\begin{tabular}{@{}p{2.0cm}p{3.5cm}p{3.5cm}p{3.7cm}@{}}
\toprule
& Whole query (one model for the request) & Repeated steps (same kind of work each step) & Distinct stages (different sub-tasks)\\
\midrule
Per instance, learned & RouteLLM (ICLR 25); AutoMix (NeurIPS 24, cascade); Router-R1 (NeurIPS 25, LLM router, RL, on HotpotQA); Conformal LLM Routing (ACL 26 SRW; PAC guarantee, dataset-level, question-only) & TRIM (2026; math reasoning; decides after the cheap model writes each step; step grader; no risk guarantee in the objective); R2-Reasoner, DoT (2025; sub-tasks a decomposer invents per query, not fixed stages) & \textbf{This work.} Nearest misses: MasRouter (ACL 25; decides upfront, never reads stage output); TRAQ (NAACL 24; per-stage conformal sets, no model choice)\\
\addlinespace
Per dataset, fixed & One model per benchmark (standard practice) & RSD, SpecReason (ICML / NeurIPS 25; accept each cheap step by a fixed threshold rule, not learned) & AgentTTS (NeurIPS 25; search; one setting per dataset, incl.\ HotpotQA); LLMSelector, Cognify (static per-module allocation)\\
\bottomrule
\end{tabular}
\end{table}

\section{Risk rubric}
\label{app:rubric}
Each instance is scored 1--5 on three items (5 = AI-ready, low risk; 1 = human-led, high risk). \textbf{$Q_1$ Impact of error}: 5 low-stakes (minor inconvenience even if errors spread); 4 mostly low-stakes; 3 moderate (meaningful per-instance impact or significant cascade); 2 high (significant impact or severe cascade); 1 critical (severe impact and/or catastrophic cascade). \textbf{$Q_2$ Detectability}: 5 errors surface immediately; 4 visible with light validation; 3 detectable with deliberate review; 2 require extensive validation; 1 invisible until significant damage. \textbf{$Q_3$ Reversibility} (new for routing): 5 undone immediately with no residual effect; 4 undone with minor effort; 3 some effects persist; 2 mostly irreversible; 1 irreversible. $\tau=(15-\sum Q)/12$; floor 0.8 if $Q_1{=}1$; minimum 0.10; maximum 0.95. On MedAgentBench, medication orders (two task types) score $\tau\ge0.8$.

\section{Preliminary experiments}
\label{app:prelim}
\begin{figure}[!htb]
\centering
\begin{minipage}[t]{0.55\linewidth}\vspace{0pt}
\includegraphics[width=\linewidth]{figs/headroom.pdf}
\end{minipage}\hfill
\begin{minipage}[t]{0.42\linewidth}\vspace{2mm}
\caption{\textbf{Evidence that headroom exists.} Preliminary results on 1,481 held-out HotpotQA questions using a two-stage retrieval--answering task with Qwen2.5-1.5B-Instruct (\W) and Qwen2.5-7B-Instruct (\Sg); flat cost per model call; lenient match under the standard ruler used for all results. The oracle---the best choice per question, known only in hindsight---is an upper bound, not a method: 68.60\% accuracy against 59.55\% for always Strong$\cdot$Strong ($+9.1$) at about a third of the cost. Unwinnable questions (31.4\%) are charged at the Weak$\cdot$Weak price. Token-aware costs: Table~\ref{tab:ladder}.}
\label{fig:headroom}
\end{minipage}
\end{figure}



Setup as in \S\ref{sec:insights}. Table~\ref{tab:insights} summarises how each insight shapes the RightRoute design.

\begin{table}[h]
\caption{The three insights and how each shapes the RightRoute design.}
\label{tab:insights}
\centering\small
\begin{tabular}{@{}p{3.6cm}p{9.4cm}@{}}
\toprule
Insight & How it shapes the RightRoute design\\
\midrule
1.\ Strong is not always best & \textbf{Outcome reward:} the router is trained on a reward that prices accuracy, cost and risk, not to imitate the strongest model, because the right executor changes by instance and stage.\\
2.\ Input decides learnability & \textbf{Weak first, one router per stage:} each router reads the Weak model's attempt, never the raw input, and asks whether it is enough.\\
3.\ Upstream choice sets downstream accuracy and cost & \textbf{Sequential, branch-priced decisions:} stages are decided in order, and every cost in the reward is conditional on the branch taken.\\
\bottomrule
\end{tabular}
\end{table}

\textbf{Worked example.} Table~\ref{tab:example} shows why per-instance choice helps on five illustrative questions. The best configuration differs per question; on q4 only the all-Weak path is correct, so always-Strong is not an upper bound on accuracy. The oracle answers 4/5 against 3/5 for always-Strong.

\begin{table}[h]
\caption{Illustrative instances (\checkmark{} = correct). Flat cost per instance: \W\W{} 2.0, \W\Sg{} 5.7, \Sg\W{} 5.7, \Sg\Sg{} 9.4 (weak pass 1.0, strong pass 4.7). Unwinnable instances (q5) are charged at the \W\W{} price. Example question: ``Which magazine was started first, Arthur's Magazine or First for Women?''}
\label{tab:example}
\centering\small
\begin{tabular}{@{}lccccc@{}}
\toprule
& \W\W & \W\Sg & \Sg\W & \Sg\Sg & Oracle choice\\
\midrule
q1 & \checkmark & \checkmark & \checkmark & \checkmark & \W\W\\
q2 & $\times$ & $\times$ & \checkmark & \checkmark & \Sg\W\\
q3 & $\times$ & $\times$ & $\times$ & \checkmark & \Sg\Sg\\
q4 & \checkmark & $\times$ & $\times$ & $\times$ & \W\W\\
q5 & $\times$ & $\times$ & $\times$ & $\times$ & none (charged \W\W)\\
\midrule
Correct & 2/5 & 1/5 & 2/5 & 3/5 & \textbf{4/5}\\
\bottomrule
\end{tabular}
\end{table}

\textbf{Fixed-policy ladder and cost families.} Table~\ref{tab:ladder} is the ruler every experiment is measured against. We report two cost families. The \emph{flat index} charges 1.0 per weak pass and 4.7 per strong pass (parameter ratio). The \emph{token-aware} family (${\approx}$FLOPs) is measured from the grid's recorded token counts (Table~\ref{tab:passcost}); because the answer-stage prompt carries the selected evidence and weak retrieval selects about $2.6\times$ more tokens, answer-stage prices depend on the retrieval choice. All cost claims use the token-aware family; on it the oracle costs 31.6\% of always-Strong (31.9\% flat). 31.4\% of instances are answered correctly by no configuration and are charged the cheapest configuration's price as a floor.

\begin{table}[h]
\caption{Fixed-policy ladder (HotpotQA held-out test, $n{=}1{,}481$, Qwen2.5 1.5B/7B, standard ruler).}
\label{tab:ladder}
\centering\small
\begin{tabular}{@{}lccc p{4.6cm}@{}}
\toprule
Policy & Acc.\ (\%) & Cost, flat (\%) & Cost, token-aware (\%) & Role\\
\midrule
Always \W\W & 37.68 & 21.3 & 18.3 & cheap floor\\
Always \Sg\W & 48.48 & ${\sim}61$ & 98.1 & nearly always-Strong-priced: strong retrieval alone is ${\sim}98\%$ of spend\\
Always \W\Sg & 49.29 & ${\sim}61$ & 23.1 & $+0.8$ over \Sg\W{} at a quarter of its real price\\
Always \Sg\Sg & 59.55 & 100 & 100 & the accuracy bar; match it or beat its price\\
Oracle & 68.60 & 31.9 & 31.6 & ceiling of perfect routing (hindsight)\\
\bottomrule
\end{tabular}
\end{table}

\begin{table}[h]
\caption{Cost per model pass (relative units).}
\label{tab:passcost}
\centering\small
\begin{tabular}{@{}lcc@{}}
\toprule
Pass & Flat index & Token-aware (${\approx}$FLOPs)\\
\midrule
Weak retrieval & 1.0 & 43.1\\
Strong retrieval & 4.7 & 201.3\\
Weak answering & 1.0 & 2.6 on weak evidence; 1.0 on strong evidence\\
Strong answering & 4.7 & 12.1 on weak evidence; 4.7 on strong evidence\\
\bottomrule
\end{tabular}
\end{table}

\textbf{Learnability probe.} The probe is the router architecture itself (Qwen2.5-1.5B-Instruct, LoRA rank 16) trained on the coverage label (does the Weak retriever's selection contain a gold answer?) from the Weak retriever's selected passages on the 4,814 training questions, and evaluated on the 1,110 validation questions (no cross-validation; checkpoint by validation AUC): AUC 0.756 at an 8,192-token call budget and 0.766 at 16,384; the trained retrieval router then reached 0.771 on validation. By contrast, routers shown the raw question and 100-passage pool and asked which model would win collapsed to a constant policy (AUC ${\le}0.58$), and six cross-validated feature-family probes on all 7,405 instances (question embeddings, BM25 pool statistics, Weak-model behaviour, their combination, a closed-book probe inside the Strong model, and self-consistency) predicted which model wins at 0.51--0.57 while predicting difficulty at 0.61--0.69. This motivated the compact router input and the Weak-first design (\S\ref{sec:insights}).


\textbf{Scorer history.} Before 13 Oct 2026 these results were scored with an earlier, stricter grid rule (the version shown in our proposal): \W\W{} 31.80, \Sg\W{} 41.73, \W\Sg{} 43.55, \Sg\Sg{} 53.54, oracle 61.24 at 34.5\% cost, with 38.8\% unwinnable. Re-grading the same stored answers under the standard ruler, which also accepts a partial answer contained in the gold, moved almost every number up by 6--7 points without changing the ranking or the headroom (about $+8$--$9$ points). The grid rule was replayed over all 29,620 stored verdicts with 0 mismatches and is kept as an audit key.


\section{Additional results}
\label{app:extra}
\begin{table}[h]
\caption{\textbf{The three claims, the evidence for each, and its main caveat} (held-out tests).}
\label{tab:claims}
\centering\small
\begin{tabular}{@{}p{3.3cm}p{5.3cm}p{4.4cm}@{}}
\toprule
Claim & Evidence & Caveat\\
\midrule
1.\ Reaches the always-Strong bar at lower cost; above it on MedAgentBench & GRPO(Reward) router: MedAgentBench 61.3 vs.\ 54.0\% at 66.6\% cost ($+7.3$); HotpotQA 75.89 vs.\ 76.8\% at 69.7\% cost\iftwowiki; 2Wiki 63.0 vs.\ 65.8\% at 85\% (bar not reached)\fi & No significance test: the MedAgentBench gain ($+7.3$) is within the $\pm7.7$-point interval at $n{=}150$, so it is read as consistent, not proven; ``reaches'' on HotpotQA is a tie within $\pm2.2$ points, not a formal non-inferiority result\\
2.\ Risk-gated W/S/H with human escalation & Priced router: all 30 high-risk orders reached a human, 0 model failures (vs.\ 21 without the risk term); identical across 3 seeds & Human escalation counted as correct; 9 of 30 escalations unnecessary; human and failure prices assumed; 30 instances, one benchmark\\
3.\ Three choices, one per objective & Accuracy: fine-tuned executors 81.3\% (HotpotQA), 76.7\% model-only and 97.3\%$^\ast$ with a human (MedAgentBench). Cost: router at 54--79\% of always-Strong. Deployability: ${\approx}4$M trained parameters, ${\sim}2$\,h on one GPU & Each choice is compared with always-Strong on the same held-out instances; published figures are context only (different test sets)\iftwowiki; on 2Wiki no approach beats the bar\fi\\
\bottomrule
\end{tabular}
\end{table}

\textbf{Full router-size grid.} Table~\ref{tab:size} extends Table~\ref{tab:routersize} to all 16 pairings. Decision latency per router call: 11.6 / 18.4 / 29.6 / 66.4\,ms (retrieval) and 11.7 / 19.3 / 33.2 / 72.6\,ms (answering) for 0.5B / 1.5B / 3B / 7B.

\begin{table}[h]
\caption{Router size ablation (validation, SFT (label) recipe): cost to reach always-Strong accuracy, \% of always-Strong (929 units), both router calls charged (parameter-scaled from the 7B pass = 77 units). Accuracy spread across all 16 cells: 76.76--77.03\%.}
\label{tab:size}
\centering\small
\begin{tabular}{@{}lcccc@{}}
\toprule
Retrieval router $\downarrow$ / Answering router $\rightarrow$ & 0.5B & 1.5B & 3B & 7B\\
\midrule
7B & 115.2 & 116.2 & 118.1 & 122.6\\
3B & 93.6 & 92.1 & 96.7 & 101.5\\
1.5B & 86.9 & \textbf{74.3} & 89.8 & 94.6\\
0.5B & 84.4 & 84.6 & 87.3 & 92.1\\
\bottomrule
\end{tabular}
\end{table}

\begin{figure}[h]
\centering
\includegraphics[width=0.9\linewidth]{figs/ft_exec.pdf}
\caption{Router-free ablation on HotpotQA (held-out test, $n{=}1{,}481$; accuracy = lenient match; cost = measured compute as \% of always-Strong; single seed). Learned escalation rules (79.6\%) never beat the best fixed fine-tuned pair (81.2\%), and very small models on the escalation ladder kept cost low but lost accuracy (67.4\%). AgentTTS's passage file places gold passages first, worth $+2.6$ points to untrained models and none to ours.}
\label{fig:ftexec}
\end{figure}

\label{app:systemone-results}\textbf{System One judges.} Tables~\ref{tab:judges-hotpot} and~\ref{tab:judges-mab} give the full judge results. Besides operating points, judges are compared by AIQ \citep{hu2024routerbench}: the area under the accuracy--cost curve traced by a router's threshold, normalised over the cost range from always-Weak to always-Strong (weak-first prices on HotpotQA, 80 to 1,009 units), with 95\% paired-bootstrap intervals on differences.

\textbf{Judges vs.\ routers.} On the main executors, before their own compute is charged, the fine-tuned E4B judge that reads the Weak answer beats both LLM routers (0.736 vs.\ 0.726 for SFT (label) + Reward, $[+0.003, +0.019]$, and 0.724 for SFT (label), $[+0.004, +0.021]$), and the 2B judge beats the label router ($[+0.001, +0.014]$) and ties the reward router. With their price (three passes: 20.2 and 4.7 units a question) both fall to 0.730 and tie both routers, whose own compute is not charged; a retrain of the E4B judge for the published weights narrowly beats both with its price (label router $[+0.001, +0.016]$, reward router $[+0.0002, +0.014]$), so the two draws support ``ties or narrowly beats''. On AIQ the judges rank questions as well as the routers, but at the routers' operating rules only the E4B ensemble finds a cheaper point at the label router's accuracy (76.70\% at 737.5 vs.\ 834.8 units; Table~\ref{tab:judges-hotpot}). nano-jev trails both routers (label $[-0.011, -0.0004]$, reward $[-0.013, -0.002]$) at $1/38$ to $1/75$ of a Decider's compute per decision. Its case is deployment: on 8 CPU threads it decides in 123\,ms against 4.8\,s for the 1.5B router, whereas on an L40S the gap is small (6.4 vs.\ 11.5\,ms), and the Deciders need a GPU.

\begin{table}[h]
\caption{\textbf{System One judges on HotpotQA} (held-out test, $n{=}1{,}481$; higher AIQ is better). ``With price'' adds the judge's own compute at measured tokens (nano-jev prices at 0). \emph{Main executors:} judges are three-seed ensembles (three passes priced); the operating point is chosen on validation with the LLM routers' own rules (peak for the label router, parity for the reward router; for the judges both rules give the same point) and read once on test, excluding the judge's passes. \emph{Preliminary pair} (Qwen2.5-1.5B on its top two passages, Qwen2.5-7B on all ten): AIQ of the three-seed mean; with price, one fine-tuned judge (range over seeds), since three passes never pay on this pair.}
\label{tab:judges-hotpot}
\centering\small
\begin{tabular}{@{}llccc@{}}
\toprule
Executors & Router & AIQ & With price & Operating point (acc.\ at units)\\
\midrule
\multirow{8}{*}{Main} & Decider E4B, qpa & \textbf{0.736} & 0.730 & 76.70 at 737.5\\
& Decider 2B, qpa & 0.731 & 0.730 & all-\Sg{} (76.77 at 1,009)\\
& Decider E4B / 2B, qp & 0.723 / 0.722 & 0.717 / 0.720 & all-\Sg\\
& nano-jev, qpa / qp & 0.719 / 0.718 & 0.719 / 0.718 & all-\Sg\\
& LLM router, SFT (label) + Reward & 0.726 & uncharged & 75.96 at 664.7\\
& LLM router, SFT (label) & 0.724 & uncharged & 76.70 at 834.8\\
& Fixed paths (hull) & 0.693 & -- & --\\
& Always-Strong & -- & -- & 76.77 at 929\\
\midrule
\multirow{9}{*}{Preliminary} & Decider E4B fine-tuned, qpa & \textbf{0.711} & 0.654--0.655 & --\\
& Decider 2B fine-tuned, qpa & 0.696 & \textbf{0.677--0.682} & --\\
& Decider E4B zero-shot & 0.696 & 0.650 & --\\
& nano-jev, qpa / qp & 0.677 / 0.674 & 0.677 / 0.674 & --\\
& nano-jev, qa / q (no passages) & 0.646 / 0.641 & -- & --\\
& LR cascade & 0.676 & -- & --\\
& Entropy threshold & 0.665 & -- & --\\
& MS MARCO cross-encoder, qp & 0.659 & -- & --\\
& Random & 0.638 & -- & --\\
\bottomrule
\end{tabular}
\end{table}

\textbf{Preliminary pair.} Before its price the fine-tuned E4B judge is the best router (0.711, $+0.02$ to $+0.06$ over entropy, the LR cascade and nano-jev), but one E4B pass costs 4.4 units, nearly half the 7B's answer, and drops it to 0.655. With its price, one fine-tuned 2B qpa judge (1.3 units a call) is the best deployable router (0.677--0.682): above entropy for all three seeds (weakest $[+0.001, +0.024]$) and level with the LR cascade. nano-jev ties entropy and the LR cascade at no price. Judges without passages (q, qa) sit near random, far below qp: the judge-level counterpart of Insight~2. On the full 7,405-question development set with GFLOPs costs, the nano-jev qp judge (0.673) beat entropy (0.662; $[+0.006, +0.016]$) and the matched pre-generation LR router (0.665; $[+0.004, +0.011]$), and tied the best cascade LR router (0.672) at 3\% of the small model's compute.

\begin{table}[h]
\caption{\textbf{System One judges on MedAgentBench, and judge cost.} \emph{Top:} held-out test ($n{=}150$, 49 patients); point A is the most accurate validation threshold (cheapest on ties), committed before the single test read. Cost = \% of always-escalate under weak-first accounting (deliver = Weak episode; escalate = Weak + Strong episodes); $\times1.23$ gives the Strong-alone basis of Table~\ref{tab:main} (e.g.\ 59.0 $\to$ 72.6). AIQ is over the weak-first cost range (random 0.330). Fine-tuned judges are seed ensembles (five seeds for the 33M models, three for the Deciders) trained on 90 tasks; their own price changes AIQ by at most 0.002. Pure-T0: the Strong model rerun without repetition penalty (always-escalate 56.67 vs.\ 54.00 on these tasks), judges retrained, a second read of the same tasks reported as sensitivity. \emph{Bottom:} GFLOPs for one pass on the preliminary pair (qp input, $2\times$ parameters $\times$ tokens); latency p50 per decision on real test inputs; price for three passes on the main executors.}
\label{tab:judges-mab}
\centering\small
\resizebox{\linewidth}{!}{%
\begin{tabular}{@{}lcccccc@{}}
\toprule
& \multicolumn{3}{c}{Canonical table} & \multicolumn{3}{c}{Pure-T0 table}\\
\cmidrule(lr){2-4}\cmidrule(lr){5-7}
Router (point A) & Acc.\ (\%) & Cost (\%) & AIQ & Acc.\ (\%) & Cost (\%) & AIQ\\
\midrule
Always-escalate & 54.00 & 100 & -- & 56.67 & 100 & --\\
Category $\times$ finished (template lookup) & 63.33 & 67.7 & 0.529 & 66.00 & 67.7 & 0.543\\
nano-jev, fine-tuned & 63.33 & 59.0 & 0.531 & 66.00 & 60.9 & 0.547\\
Decider 2B, fine-tuned & 62.00 & 58.1 & 0.532 & 66.00 & 59.7 & 0.545\\
Decider E4B, fine-tuned & 62.67 & 56.3 & 0.532 & 66.00 & 62.1 & 0.545\\
MiniLM, fine-tuned (no decision pretraining) & 59.33 & 52.1 & 0.521 & 65.33 & 59.1 & 0.544\\
MS MARCO MiniLM, fine-tuned & 62.67 & 85.5 & 0.522 & 64.67 & 87.4 & 0.544\\
nano-jev, zero-shot & 53.33 & 99.6 & 0.383 & 56.00 & 99.6 & 0.398\\
Decider 2B, zero-shot & 50.67 & 97.1 & 0.400 & 53.33 & 97.1 & 0.418\\
Decider E4B, zero-shot & 53.33 & 99.6 & 0.431 & 56.00 & 99.6 & 0.451\\
\bottomrule
\end{tabular}}

\vspace{2mm}
\begin{tabular}{@{}lccccc@{}}
\toprule
Judge & Parameters (trained) & GFLOPs & L40S, batch 1 & 8 CPU threads & Price, menu A\\
\midrule
nano-jev & 33.4M (all) & 32 & 6.4\,ms & 123\,ms & 0\\
Strands Decider 2B & 2B (17.9M) & 1,215 & 38.6\,ms & needs a GPU & 4.7\\
Strands Decider E4B & 4B active (109M) & 2,401 & 47.7\,ms & needs a GPU & 20.2\\
LLM router (Qwen2.5-1.5B) & 1.5B (4.4M) & -- & 11.5\,ms & 4,787\,ms & uncharged\\
\bottomrule
\end{tabular}
\end{table}

\textbf{MedAgentBench.} The fine-tuned judges reach 62.0--63.3\% (oracle 63.3\%) at 56--59\% of always-escalate cost. They match the template lookup that escalates by task type and by whether the Weak agent finished (Deciders $+0.003$, $[+0.002, +0.007]$; nano-jev $+0.003$, $[-0.0005, +0.007]$), and nano-jev beats the MS MARCO start ($+0.009$, $[+0.001, +0.016]$). Most of the signal is at the task-type level: the Weak model is right almost only on one task type and partly on another. Paired on the same tasks with the GRPO(Reward) router (61.33\% at 54.1\% of always-escalate), fine-tuned nano-jev is $+2.0$ points ($[0.00, +4.52]$) at $+4.9$ points of cost ($[+2.3, +7.9]$), so no superiority is claimed. Zero-shot, the E4B closes 50\% of the AIQ gap between random and fine-tuned nano-jev, and nano-jev 26\%. The pure-temperature rerun lifts every router by 2--6 points; the fine-tuned judges still match or edge the template lookup and the zero-shot judges still trail it, with differences of 0.001--0.005 AIQ, so these results are exploratory.

\textbf{What the typed decisions add.} A pre-registered study on the preliminary pair (HotpotQA exploratory; MuSiQue \citep{trivedi2022musique} and 2WikiMultiHopQA \citep{ho2020constructing} confirmatory; margin 0.010 AIQ, Holm correction within each family) tested which properties of the judge hold beyond HotpotQA. Confirmed on both new datasets: (i) starting from nano-jev beats the same 33M MiniLM without decision training (area under the AIQ-vs-label-budget curve, qp: $+0.021$ / $+0.010$ / $+0.019$ on HotpotQA / MuSiQue / 2Wiki), and nano-jev also beats the MS MARCO start \citep{bajaj2016msmarco} in every cell; (ii) the passages, not the question, carry the escalation signal (qp $-$ q: $+0.011$ / $+0.025$ on MuSiQue / 2Wiki); (iii) the text-only judge is non-inferior to entropy, the LR routers and every cheap signal from the Weak model, including P(True), an NLI check and a probe on its hidden states (closest, the probe: $+0.009$ / $+0.005$ / $+0.008$), and with their extra calls priced, self-consistency and semantic entropy fall below random; (iv) without routing labels, one untrained typed question recovers 77\% / 86\% / 83\% of a trained judge's gain over random. Not supported: a larger advantage with few labels, a benefit from the typed head itself, and useful distribution-free risk control (Appendix~\ref{app:negative}).

\textbf{Contamination.} nano-jev's pretraining saw 1,000 HotpotQA development questions, 178 of them in our HotpotQA test; on the other 1,303 every comparison keeps its order (main executors: E4B qpa 0.739, 2B qpa 0.735, LLM routers 0.727 / 0.729). MuSiQue appears only in nano-jev's held-out evaluation; the Deciders' training data includes MuSiQue, so no Decider number enters a confirmatory test.



\textbf{Routing-score diagnostics (MedAgentBench GRPO routers).} AUC 0.93 / 1.00 / 0.91 and calibration error 0.34 / 0.03--0.04 for the full-attempt, risk-term and second-round runs; decision latency 19\,ms median, 25\,ms p95 (1.5B router). Probabilities bunch near 0 and 1.

\section{Metrics}
\label{app:metrics}
\textbf{Notes to Table~\ref{tab:main}.} Cost is task-model compute as \% of the Strong model alone on the same instances (HotpotQA: 929 units per question; MedAgentBench: the Strong agent's cost per task); the routers' and judges' own compute is reported below. $^\ddagger$AgentTTS Table~5 prices: 77 + 3 = 80 units. $^\S$Published on a since-retired hosted FP8 serving (Together); our local AWQ serving of the same model scores 56.7 (all 300) and 54.0 (our test). $^\P$Model cost as \% of Strong alone, plus human time in money (SGD 183.75 for the 30 high-risk tasks escalated to a human); compute and human cost are kept in separate units, with no conversion between them. The AgentTTS (reproduced) row is on a 500-question slice of the same dev set (AgentTTS's own slice is unreleased), not our held-out test. Fine-tuned executors: Qwen 7B retriever $\to$ Llama 70B answerer (HotpotQA); 7B alone (MedAgentBench). MedAgentBench costs from teammates' pipelines were reported against always-escalate (Weak attempt + Strong) and are converted to the Strong-alone basis by $\times1.23$ (1,566.6 / 1,273.8 units on the held-out 150).

Composed accuracy (lenient headline; exact match supplementary); cost per instance; cost including decision overhead (router calls, judge probe; $\approx0$ for classifiers); cost to reach always-Strong accuracy (``not reached'' if no operating point attains it); calibration; AUC (a training signal, not a final result); decision latency. Metrics 1--2 decide which approach wins (accuracy is the constraint, cost the score); 3--4 explain behaviour. \emph{Router compute.} Headline costs count task-model compute only. The LLM router's own compute adds about 3.6 points of always-Strong on HotpotQA (both router calls) and takes about 19\,ms per decision on MedAgentBench. The System One E4B judge's own three passes add 20.2 units a question on HotpotQA (2.2 points; 81.6\% with them), and its HotpotQA point needs the three-seed ensemble (no single seed reaches it); nano-jev's own price rounds to zero on both benchmarks (Appendix~\ref{app:systemone-results}).

\section{Implementation details}
\label{app:impl}
\label{app:families}\textbf{Alternative families: full descriptions} (summarised in \S\ref{sec:families}).

\textbf{Non-LLM router.} A lightweight classifier makes the same per-stage choice from features available at decision time: lexical (TF-IDF n-grams), structural (length, vocabulary, formatting statistics) and stage-aware (execution progress, retrieval choice). On HotpotQA we train linear and MLP classifiers with SFT (label), SFT + reward and GRPO; on MedAgentBench we additionally fit random forest, SVM and logistic-regression classifiers by reward-weighted supervised imitation (not GRPO). This is the baseline that tests whether an LLM is needed at all.

\textbf{System One judge.} A small decision model replaces the router: it scores a typed question's options over a text state, without generating. At each stage it asks \emph{will the Weak result do?} and escalates if $P(\text{escalate})\ge\theta$, reading the question and the Weak retriever's passages (retrieval), the chosen branch's passages and optionally the Weak answer (answer; qp or qpa), or the task and the Weak agent's attempt (MedAgentBench). We use nano-jev \citep{nanojev2026}, a 33.4M-parameter MiniLM cross-encoder \citep{wang2020minilm} pretrained on typed decisions over HotpotQA, SQuAD 2.0 and MultiNLI, and the 2B and E4B Strands Deciders \citep{strandsdecider2026}, LoRA-adapted LLM torsos with a pointer head over options. Judges train on the same per-stage labels as SFT (label), are averaged over seeds and pay for their own compute (Appendices~\ref{app:extra},~\ref{app:systemone}).

\textbf{Fine-tuned executors with learned escalation.} \andy{Instead of routing between fixed models, we fine-tune the executors themselves and let rules decide when to escalate. On HotpotQA we fine-tune three retrievers (Qwen2.5 7B/32B/72B) and three answerers (Llama 3.2 3B, Llama 3.1 8B/70B) on questions from the official training set: each retriever learns to pick, from the 100 candidates, the two passages that HotpotQA marks as holding the supporting facts, and each answerer learns to give the dataset's reference answer (LoRA on 12,000 questions; QLoRA on 2,000 for the 32B--72B models). At test time the smallest models go first, and rules decide after each step whether to add a larger one, possibly skipping a size. The retrieved shortlists are merged (unanimous, else majority, else union); an answer that agrees with an earlier one is final, and otherwise a majority vote decides, with ties going to the larger model. Rules are fixed or learnable: an LLM planner (Qwen2.5-72B) proposes up to three rule changes per round over four rounds, and a change is kept only if it improves both random halves of the validation set, respects guardrails (e.g., at most 60\% of traffic to the largest model) and beats random escalation of the same volume. MedAgentBench follows the same design with one agent model per rung and a human as the last rung on high-risk tasks; Table~\ref{tab:ftexec} (Appendix~\ref{app:rules}) sets out both setups.} \andynote{rewritten to focus on HotpotQA and to say what the models learn; MedAgentBench details moved to Table~\ref{tab:ftexec}}

\textbf{Reward settings.} HotpotQA: $Q$ = token-F1 (lenient match remains the evaluation ruler); $\lambda{=}0.10$ (selections identical across 0.05--0.11; $\ge0.12$ drops below always-Strong); $C_{SS}{=}929$; $\lambda_t{=}0$; weak-first accounting (an escalation costs the Weak attempt plus the Strong call; Strong everywhere = 1,009 units). MedAgentBench: $Q$ = official pass/fail; $\lambda{=}0.10$; per-instance $C_{SS,i}$; weak-first accounting (deliver $= C_\W$, escalate $= C_\W + C_\Sg$); $\tau_{\text{gate}}{=}0.8$; $\lambda_t{=}1$.

\textbf{LLM router: prompts and training hyperparameters.}

\emph{Router calls.} Every router call is a system prompt plus one stage instruction, rendered through the Qwen chat template up to the opened assistant turn. The router generates nothing: the next-token logits for \texttt{W} and \texttt{S} are read and softmaxed over those two tokens. \texttt{\{weak\_name\}} and \texttt{\{strong\_name\}} are the actual model names (on HotpotQA with the main executors: Qwen2.5-7B / Qwen2.5-72B at retrieval, Llama-3.2-3B / Llama-3.1-70B at answering).

\emph{System prompts (HotpotQA).}
\begin{quote}\small
\textbf{SFT router prompt} (SFT (label) and SFT (label) + Reward): \texttt{You route a question-answering pipeline that runs the weak model first and escalates only when needed. Each stage runs either \{weak\_name\} (weak, cheap) or \{strong\_name\} (strong, expensive). From the weak model's output shown below, choose the cheapest option that still lets the pipeline answer correctly.}

\textbf{Reward-objective prompt} (GRPO(Reward)): \texttt{You route a question-answering pipeline. Each stage runs either \{weak\_name\} or \{strong\_name\}. Your objective: the pipeline should answer correctly while spending as little compute as possible - correctness first, cost second. Based on the information shown below, choose one option.}

\textbf{Three-insights variant} (ablation; the SFT router prompt followed by): \texttt{Routing principles: (1) Begin with the weak model first. (2) Escalate to the strong model only when the weak output is unlikely to be correct. (3) If even the strong model would be unlikely to answer correctly, keep the weak choice - escalation that cannot fix the answer only adds cost.}
\end{quote}

\emph{Stage instructions (HotpotQA).}
\begin{quote}\small\ttfamily
Stage: RETRIEVAL REVIEW --- the weak retriever has already selected references from the 100-passage pool.\\
W = keep the weak selections (no further retrieval spend)\\
S = escalate: run strong retrieval to reselect\\
Question: \{q\}\\
Weak retriever selected \{n\_sel\} of 100 passages: \{selections\}\\
Answer with one letter only.

\medskip
Stage: ANSWER REVIEW --- the weak answerer has already produced an answer from the references below.\\
W = deliver the weak answer as final\\
S = escalate: run strong answering\\
Question: \{q\}\\
References in use: \{evidence\}\\
Weak answer: \{weak\_answer\}\\
Answer with one letter only.
\end{quote}
\texttt{\{selections\}} lists the Weak retriever's passages as \texttt{[n] text}; \texttt{\{evidence\}} is the reference block for the branch taken. Each passage shown to a router is capped at 600 characters (97.7\% of gold answers survive the cap, chosen by audit over all 7,405 instances); an assembled router call is capped at 16,384 tokens, left-truncated (no call reached it; the longest was 15,261).

\emph{MedAgentBench prompt.} System prompt (\texttt{\{weak\_name\}} = Qwen2.5-7B, \texttt{\{strong\_name\}} = Qwen2.5-72B-AWQ):
\begin{quote}\small\ttfamily
You route a medical record-system task. The weak model \{weak\_name\} has already attempted it; choose whether to deliver its result or escalate to \{strong\_name\}. Your objective: the task should be completed correctly while spending as little compute as possible - correctness first, cost second. Based on the information shown below, choose one option.
\end{quote}
The user message is the task instruction (plus \texttt{Context: \ldots} when the task has one), a blank line, the view of the Weak attempt, then \texttt{Options: W) Deliver the weak model's result / S) Escalate to the strong model / Answer with a single letter: W or S.} The view is either \texttt{The weak model finished after \{r\} round(s). Its final answer: \{answer\}} or \texttt{The weak model did NOT finish: it stopped after \{r\} round(s) without calling FINISH. Its last action: \{action\}}, with the answer or action cut at 400 characters. The second-round-view router builds the same view from the attempt cut at its second round. The router never sees the grader's verdict. The risk-term routers add a third option on $\tau\ge0.8$ instances only, worded in the prompt as \texttt{H) Hand the task to a person}.

\emph{Router model and adapters.} Qwen2.5-1.5B-Instruct with the base frozen; LoRA rank 16, $\alpha{=}32$, dropout 0.05, on \texttt{q\_proj}, \texttt{k\_proj}, \texttt{v\_proj} and \texttt{o\_proj}: 4,358,144 trainable parameters (0.28\%). One adapter per stage on HotpotQA; one router on MedAgentBench.

\begin{table}[!htb]
\caption{\textbf{LLM router training hyperparameters.} Executors decode greedily ($K{=}1$); the $K{>}1$ cells sample at $T{=}0.7$, top-$p$ 0.9 (temperature ablation $T\in\{0.3, 1.0, 1.3\}$). The router reads logits only and generates nothing.}
\label{tab:hparams}
\centering\scriptsize
\resizebox{\linewidth}{!}{%
\begin{tabular}{@{}p{2.3cm}p{2.9cm}p{2.3cm}p{3.2cm}p{3.2cm}@{}}
\toprule
& SFT (label) & SFT (label) + Reward & GRPO(Reward), HotpotQA & GRPO(Reward), MedAgentBench\\
\midrule
Objective & class-weighted cross-entropy on per-stage labels (weights $N/(2\cdot\text{count})$; X trained as \W) & expected reward $R$ over the cached grid & GRPO on Eq.~\ref{eq:reward} & GRPO on Eq.~\ref{eq:reward}\\
Initialisation & fresh LoRA & from the SFT (label) adapters & fresh LoRA (label-free) & fresh LoRA (label-free)\\
Rollouts per group & -- & -- & 8 per question, decisions sampled at both stages ($2\times2$ paths) & 8 per task\\
Advantage & -- & -- & $R$ minus group mean (no std normalisation) & same\\
Entropy bonus $\beta$ & -- & 0.03 & 0.03 & 0.03\\
Learning rate & $1\times10^{-4}$ & $1\times10^{-5}$ & $1\times10^{-5}$ & $1\times10^{-5}$\\
Grad.\ accumulation / clip & 8 / -- & 8 / 1.0 & 8 / 1.0 & 8 / 1.0\\
Epochs & 3 & 3 & 6 (selected) & 40 (priced-risk router: 40, the first 10 with exploration)\\
Seed & 17 & 17 & 17 & 17\\
Reward constants & -- & $\lambda{=}0.10$, $C_{SS}{=}929$ & $\lambda{=}0.10$, $C_{SS}{=}929$, hinge guard & $\lambda{=}0.10$, per-task $C_{SS,i}$, hinge; priced-risk router: $\lambda_t{=}1$, $\rho{=}C_F/C_H$\\
Quality $Q$ & -- & token-F1 & token-F1 & official pass/fail grader; $Q{=}1$ for human escalation\\
Checkpoint / threshold & best validation & best validation & best validation mean $R$; $\theta_{\text{ret}}$, $\theta_{\text{qa}}$ swept on validation & best validation mean $R$; $\theta$ on validation\\
\bottomrule
\end{tabular}}
\end{table}

\textbf{Non-LLM router.} \emph{Features (HotpotQA).} Per stage, train-only TF-IDF over word 1--2-grams (min\_df 2, max\_df 0.98, 12,000 features, sublinear tf) plus six standardised text statistics (character count, token count, vocabulary size, question-word overlap, newline count, digit fraction). The retrieval view is the question plus the Weak retriever's selected passages; the answer view is the question, the selected branch's evidence and the Weak answer, plus a retrieval-branch indicator. \emph{Features (MedAgentBench).} TF-IDF 1--2-grams (at most 6,000) plus ten statistics of the same task view the LLM router sees; no outcome, cost, risk or Strong-output field enters the features. \emph{Architectures.} Linear softmax, MLP (64) and MLP (128$\to$64) with ReLU; the main HotpotQA row uses MLP (128$\to$64), about 1.55M parameters per stage (3.1M for two stages). On MedAgentBench we also fit random forest and extra trees (300 trees, depth $\in\{3,6,\text{none}\}$), decision tree, linear and RBF SVM ($C\in\{0.1,1,10\}$), logistic regression and multinomial naive Bayes, with hyperparameters chosen on validation reward and SVM margins calibrated by patient-grouped 3-fold logistic calibration; the reported row is a linear SVM (about 6K parameters). \emph{SFT (label).} Class-weighted cross-entropy, 3 epochs, AdamW, learning rate $10^{-4}$, batch 8, checkpoint chosen by validation AUC; the label is ``Weak is enough'' per stage, as for the LLM router. On MedAgentBench the classical families use cost-sensitive reward-optimal action imitation (sample weight $|R_W-R_S|$), which is supervised and not GRPO. \emph{GRPO.} On HotpotQA: group size 8, advantage normalised by the group mean and standard deviation, KL to a frozen reference 0.02, no PPO clipping, 2 epochs, learning rate $10^{-4}$, batch 16, cold or SFT-warm start, reward token-F1 $-0.1\,C/929$. On MedAgentBench: fresh start, group size 8, advantage = reward $-$ group mean (no standard-deviation normalisation, no KL), entropy 0.03, AdamW learning rate $10^{-5}$, 40 epochs, best checkpoint by validation reward checked every 10 updates. \emph{Selection.} Three seeds (17, 18, 19) are averaged before thresholds are chosen on validation (three stage thresholds on a 0.05 grid for HotpotQA; one threshold on a 0.01 grid for MedAgentBench); the test set is read once, on frozen executor outputs.

\label{app:systemone}\textbf{System One judges.} nano-jev v1.0 (\texttt{sdmlai/nano-jev}): 12-layer BERT-style cross-encoder, hidden 384, 12 heads, 33.4M parameters, all fine-tuned; one logit per (question+option, state) pair, truncated to 512 tokens, softmax over options; AdamW lr $5\times10^{-5}$, weight decay 0.01, 6\% warmup, batch 16, up to 3 epochs, best epoch by held-out NLL. Strands Decider: Qwen3.5-2B-Base (24 layers) or Gemma-4-E4B-it (42 layers) torso, frozen in bf16, LoRA on attention and MLP projections ($r{=}16$, 16.8M parameters; $r{=}48$, 104.6M) and an fp32 pointer head (1.05M; 3.9M); lr $10^{-4}$, weight decay 0.01, 6\% warmup, 2 epochs, batch 16, 512 tokens. \emph{HotpotQA (main executors):} one judge learns the three stage questions (retrieval; answer after Weak retrieval; answer after Strong retrieval) from the LLM router's per-stage labels on the 4,814 training questions, with 10\% held out by question to pick the epoch; thresholds come from the same 9,261-point validation sweep and peak/parity rules as the LLM routers; three seeds, scores averaged. \emph{MedAgentBench:} 90 training tasks, labelled escalate when the Weak attempt fails and the Strong one passes; epoch by validation NLL, five seeds for the 33M judges and three for the Deciders; 95\% intervals from a patient-level bootstrap (5,000 resamples). Zero-shot judges answer a typed question with the checkpoint's own calibrated probabilities (preliminary pair: \emph{is the claim grounded?}; MedAgentBench: \emph{is the proposed answer correct?}). \emph{Published judges.} The fine-tuned judges are public on Hugging Face, three seeds each: \texttt{mohameddhameem/systemone-jev-\{qp,qpa\}} (nano-jev) and \texttt{mohameddhameem/systemone-decider-\{2b,e4b\}-\{qp,qpa\}} (Deciders; \texttt{-menu-a} for the main executors), with code at \url{https://github.com/mohameddhameem/systemone-decider}; loading the published Decider folders reproduces every saved test score.

\label{app:rules}\textbf{Learned escalation rules.} Fixed retriever rule: confidence in lowest 10\% $\rightarrow$ escalate. Learnable starting rules: lowest 20\% $\rightarrow$ escalate skipping a size; shortlists disagree and the larger retriever is unsure $\rightarrow$ escalate again; answerer: empty or ``cannot be determined'' $\rightarrow$ escalate (fixed); yes/no $\rightarrow$ keep; answer words in passages $<97\%$ $\rightarrow$ escalate; confidence in lowest 15\% $\rightarrow$ escalate. Retriever ensemble: unanimous shortlist, else strict majority, else union (plain union on MedAgentBench). Change types: amend, add, remove, setting (remove and setting were tried but never accepted). Checks: threshold moves $\le0.05$ per round; better on both validation halves (each accuracy point may cost $\le10\%$ extra compute); guardrails ($\le60\%$ of traffic to the largest model, within 1 point of the best fixed pair of equal cost, $\le8$ learnable rules, each rule fires on $\ge2\%$); beats random escalation; only the best passing change per round is applied. \andy{The learnt rules and the full setup of the fine-tuned executors on both benchmarks are in Table~\ref{tab:ftexec}.} \andynote{replaces the TODO on the learnt rules}

\begin{table}[!htb]
\ifshowupdates\color{andycolor}\fi
\caption{\textbf{Fine-tuned executors with learned escalation: setup on both benchmarks.} Single runs (seed 0). Cost as in Table~\ref{tab:main}; on MedAgentBench human escalations are included at 1$\times$ always-Strong per escalation.}
\label{tab:ftexec}
\centering\scriptsize
\begin{tabular}{@{}p{0.15\linewidth}p{0.40\linewidth}p{0.40\linewidth}@{}}
\toprule
 & \textbf{HotpotQA} & \textbf{MedAgentBench} \\
\midrule
Executors & Retrievers Qwen2.5 7B / 32B / 72B; answerers Llama 3.2 3B, Llama 3.1 8B / 70B & One agent model per rung: Qwen2.5 7B / 32B / 72B; then a human \\
Training data & Questions from HotpotQA's official training set (not the 4,814 used by the routers) & Attempts of the untrained Qwen2.5-72B (greedy plus 4 samples) on the 90 training tasks that pass the official grader: 72 attempts covering 60 tasks. The 7B also learns from 45 demonstrations written by the authors for the three order types (training and validation tasks, no test task; the 8 that place an order count 3$\times$) \\
Target & Retriever: the two passages HotpotQA marks as holding the supporting facts, among 100 candidates. Answerer: the dataset's reference answer & Each turn of a passing attempt: the next record lookup or order (FHIR call), and the final answer \\
Method & LoRA on 12,000 questions up to 8B; QLoRA on 2,000 for 32B, 70B, 72B; checkpoint by validation loss (1,110 questions) & LoRA (7B), QLoRA (32B, 72B) \\
Ladder & Retrieval, then answering; smallest model first; a rule may skip a size & 7B $\to$ 32B $\to$ 72B $\to$ human; a human only on high-risk tasks (medication orders) \\
Combining & Shortlists: unanimous, else strict majority, else union. Answers: agreement with an earlier answer (same normalised answer, or F1 $\ge$ 0.8) is final; at the end, a vote including two sampled answers, ties to the larger model & Each escalated model adds its own record lookups to the earlier ones (union). Two layouts: every rung looks up and answers; or a lookup ladder followed by an answer ladder \\
Rule signals & Retriever confidence and shortlist agreement; answer confidence, empty or ``cannot be determined'' answers, yes/no answers, share of answer words found in the passages & Empty lookup, empty answer, ``not found'', malformed answer, several or no orders placed, answer confidence; on high-risk tasks, a doubtful attempt can go to a human \\
Rule learning & LLM planner (Qwen2.5-72B), up to 3 changes per round over 4 rounds, plus threshold steps from code; checks as above & No planner: one change at a time, the one that most improves validation accuracy net of cost (weight 0.10) without lowering either half of the 60 validation tasks \\
Learnt rules & None accepted: in both reported configurations (three sizes; ladder from 0.8B), no change passed the checks, so the starting rules are final & Every-rung layout: no order placed $\to$ 32B; on a high-risk task, no order placed $\to$ human. Lookup-then-answer layout: no order placed $\to$ the 32B writes the answer from the same lookups. Fixed in both: an unfinished 32B attempt $\to$ 72B \\
Result (test) & 79.6\% at 23\% of always-Strong (best fixed fine-tuned pair 81.3\%) & Every-rung: 97.3\% at 73\%, all 30 high-risk tasks to a human. Lookup-then-answer: 94.7\% at 51\%, 13 to a human \\
\bottomrule
\end{tabular}
\end{table}

\section{Compute disclosure}
\label{app:compute}
\textbf{Lab server.} One NVIDIA RTX 6000 Ada Generation (48\,GB; 46,068\,MiB usable), CUDA 12.9, Ubuntu, Python 3.10. Used for router training (SFT (label), the reward pass and GRPO on the 1.5B router with LoRA), the MedAgentBench router trainings and test reads, the Weak-model runs on HotpotQA (Qwen2.5-7B retrieval, Llama-3.2-3B answering, the $K{>}1$ samples and the temperature sweep), the AgentTTS fuser reproduction, and all CPU scoring and allocation. LoRA router training takes about 2 hours on this card.

\textbf{SMU research cluster} (Slurm). All Strong-model runs for the reported results used NVIDIA H200 (141\,GB) nodes: one H200 for the 7B retrieval sampling and two for the 70B/72B calls (Qwen2.5-72B retrieval, Llama-3.1-70B answering, and the AgentTTS reproduction). Other cards in the partition were not used for reported results.\todo{confirm whether any early grid run used an A100/H100 node rather than H200 (entries from October onward name H200)}

\textbf{MedAgentBench serving.} The Weak agent (Qwen2.5-7B) and the Strong agent (Qwen2.5-72B-Instruct-AWQ, 4-bit) were served with vLLM against the benchmark's FHIR container; the 72B-AWQ model fits on a 48\,GB card.\todo{confirm whether the 72B-AWQ serving ran on the lab card or an H200 node}

\textbf{System One judges.} Judge inference on one L40S at batch 1: 6.4\,ms (nano-jev; 123\,ms on 8 CPU threads), 38.6\,ms (2B), 47.7\,ms (E4B). System One judges train and are evaluated on SMU GPU-cluster L40S cards (46\,GB) from cached executor outputs (HotpotQA) or recorded episodes (MedAgentBench), so no LLM is rerun for routing; one fine-tuning seed takes 7--21\,min for the 2B Decider and 17--48\,min for the E4B, and all nano-jev judges on the main executors train in 11\,min. The cluster keeps no per-job accounting history, so these per-run times stand in for a GPU-hour total. \todo{GPU-hours per stage (grid generation, training, evaluation) and a total; Andy: fine-tuned executors' cards and hours; Ziming: non-LLM router compute}

\section{Negative findings}
\label{app:negative}
\begin{itemize}
\item A router reading the question plus the full 100-passage pool collapsed to a constant policy under every recipe (SFT with or without class weights, reward-weighted SFT, GRPO) and at 1.5B and 7B scale; six feature families predicted instance difficulty (AUC 0.61--0.69) but not which model helps (0.51--0.57).
\item A zero-shot 7B router ranked instances better than the trained ones (AUC 0.62) but was miscalibrated: no threshold beat always-Strong.
\item Selecting routers by AUC picks the 7B routers, the worst economically (Table~\ref{tab:size}).
\item Learned escalation rules never beat the best fixed fine-tuned pair; ``remove'' and ``setting'' rule changes were never accepted.
\item LLM fusion of sampled answers hurt small models.
\item System One: the typed decision head did not beat a plain classifier head on the same encoder (the gain comes from decision pretraining), and pretraining's advantage grew with labels rather than shrinking, so it is not a few-label effect.
\item System One: distribution-free risk control (Learn-then-Test thresholds \citep{angelopoulos2021ltt} bounding the error among answers kept on the Weak model) kept 0--0.5\% of questions at $\alpha\le0.2$; the bound holds only because the Weak model is wrong too often (39.5\% on HotpotQA) for any of its answers to qualify.
\item System One: at the LLM routers' operating rules, the nano-jev and 2B judge ensembles escalate everything on HotpotQA (no validation point reaches always-Strong accuracy for less) despite equal or better AIQ; two of the three single E4B seeds do the same, and only the three-seed E4B ensemble finds a cheaper point at the label router's accuracy.
\item On HotpotQA the risk screen found uniformly low $\tau$, so the risk term was never active there.
\end{itemize}

\textbf{Evaluation details (Limitation 4).}\label{app:evaldetails} We wrote 45 worked examples to teach the fine-tuned 7B model to place orders, kept only if the official grader accepted them; some come from the practice tasks later used to learn its escalation rules (none from the test set), so those rules were tuned on tasks the model had seen. The always-Strong baseline on MedAgentBench ran with a setting that discourages repeated words; without it, the model scores about 3 points higher (59.3\% on all 300 tasks), which would shrink our margins over it by about as much. In AgentTTS's HotpotQA files, answer passages tend to appear near the top of the 100 candidates, helping models that were not fine-tuned by about 2.6 points; every method in Table~\ref{tab:main} uses the same files, and our fine-tuned models score the same when the list is shuffled.

\end{document}